\section{Introducción y panorama del contenido}

\subsection{Ubicación de la clase y objetivos de aprendizaje}

En la Lecture 15, Chelsea Finn analiza a fondo el marco de ingeniería y las vías de gobernanza de la imitación y el aprendizaje por refuerzo jerárquicos. Esta sección no solo busca que entendamos por qué la estructura jerárquica es indispensable en las tareas de horizonte largo. También busca reproducir el ciclo cerrado «exploración estructurada → observación → gobernanza», que el ponente subraya una y otra vez a lo largo de la clase. La clase abarca datos, arquitectura, diseño de prompts y objetivos, y casos prácticos. Todas las notas se reorganizan según la lógica didáctica para facilitar su reproducción posterior.

\begin{knowledgebox}{Principios para construir notas con una jerarquía de dos niveles}
La clase se divide en: 1) motivación, 2) marco, 3) subobjetivos/imitación, 4) lenguaje + llamadas a LLM, 5) gobernanza de ingeniería. Cada bloque debe resaltar sus ideas centrales con `importantbox/knowledgebox/warningbox`. Las comparaciones se organizan con tablas o imágenes. Al final, un resumen sintetiza el contenido para facilitar el repaso.
\end{knowledgebox}

\subsection{Línea de tiempo y mapa de materiales}

\begin{table}[H]
\centering
\begin{tabular}{p{3.5cm}p{8cm}}
\toprule
Intervalo de tiempo & Contenido destacado y materiales \\
\midrule
00:00--00:12 & El ponente expone las razones de la jerarquía, long horizon y hierarchical induction bias (slide-01～05) \\
00:12--00:30 & Formalización de Options/SMDP y value estimation (slide-06～12) \\
00:30--00:45 & Jerarquía goal-conditioned; comparación entre HIRO y HAC (slide-13～20) \\
00:45--00:55 & Imitación jerárquica y descubrimiento de habilidades (slide-21～28) \\
00:55--01:12 & LLM + planificación con lenguaje, casos avanzados y recomendaciones de despliegue (slide-29～40) \\
01:12--01:20 & Checklists prácticas y cierre con Q\&A (slide-41～46) \\
\bottomrule
\end{tabular}
\caption{Línea de tiempo de la Lecture 15 y su correspondencia con las slides.}
\end{table}

\begin{figure}[H]
\centering
\includegraphics[width=0.85\textwidth]{slides-images/slide-01.jpg}
\caption{Slide de apertura de la clase: por qué merece atención la estructura jerárquica (slide-01).}
\end{figure}

\begin{importantbox}{Estrategia de uso de materiales}
Se citan primero las capturas de las slides. Si el contenido de una slide no basta, se complementa con los subtítulos o las notas. Cada imagen debe indicar su ruta dentro del directorio (por ejemplo, `slides-images/slide-xx.jpg`), y su caption debe incluir el timecode para poder ubicarla después.
\end{importantbox}

\subsection{Resumen de la sección}

Esta sección traza un mapa para repasar la clase a partir de su ubicación, la línea de tiempo y la guía de materiales. Todo el contenido posterior se desarrolla en este orden: motivación → modelo → práctica.
